\documentclass[11pt,a4paper]{article}
\usepackage[T1]{fontenc}
\usepackage{lmodern,amsmath,amssymb,amsthm,booktabs,microtype,hyperref}
\usepackage[margin=27mm]{geometry}
\setlength{\emergencystretch}{2em}
\newcommand{\Q}{\mathbf Q}\newcommand{\Z}{\mathbf Z}
\newcommand{\N}{\operatorname{N}}\newcommand{\Tr}{\operatorname{Tr}}
\newtheorem{theorem}{Theorem}\newtheorem{lemma}[theorem]{Lemma}
\title{Full auxiliary point charts for the general twenty-five-variable source}
\author{}\date{17 November 2026}
\begin{document}\maketitle
\begin{abstract}
The five auxiliary points of the general source may range over all
guarded affine points of the elliptic curve. The four root points
retain integer indices. The scalar valuation estimates and both
comparison precisions are preserved. Cubic auxiliary charts give a
complete twenty-five-variable specification of degree at most
$4924857429000192$, under the retained arithmetic premises. The
corresponding twenty-one-variable actual-trace source has a larger
witness space, while its simultaneous rational realization remains
unproved.
\end{abstract}

\section{The point roles}
Retain the general source and all its arithmetic premises
\cite{General,Unit,Terminal}. Thus
\[
 K=\Q(\sqrt2),\qquad E:y^2=x^3-(11-6\sqrt2)x,
 \qquad 2E(K)=\Z P,
\]
with the specified point $P$. Its four root points
$Q_i=[n_i]P$, for $1\le i\le4$, keep their original point charts.
Write $S=Q_2$ and $a=n_2$. The other roles are $H,U_1,\ldots,U_4$.
For each of these five points use
\[
 x=h_0+h_1\sqrt2,\qquad y=j,\qquad
 j^2=x^3-(11-6\sqrt2)x.
\]
Keep every original affine point, difference and rational denominator
guard. The auxiliary ordinate guard makes its cubic radicand and
both quadratic stages of its rank-four point algebra units.
The chart and all-character transfer statements of \cite{FullAux}
apply independently to these five blocks.

For some $h,j_1,\ldots,j_4\in\tfrac12\Z$ one has
\[
 2H=[2h]P,\qquad 2U_i=[2j_i]P.
\]
At every retained odd prime these labels are integral. The point
guards exclude $h=0,1$ and $j_i=0,n_i$ whenever the relevant
formal ratios occur.

\section{Local divisibilities}
\begin{lemma}\label{lem:local}
The scalar-odd, mixed-extraction and integer-residual divisibility
lemmas of the general source remain valid with the enlarged
auxiliary charts.
\end{lemma}
\begin{proof}
The scalar functions are the two rational components of each of
\[
 f(H),\qquad f(S-P),\qquad \zeta(S-P)/\zeta(H).
\]
Their valuation extraction uses only the displayed rational
functions, so is unchanged. When the first four components are
positive, $H$ and $S-P$ are in the formal group at each place
above the selected prime. The doubled group identity is
\[
 [2h](S-P)=[a-1](2H).
\]
The extended elliptic logarithm kills torsion. Since two is a
local unit and $h\ne0$, this gives
\[
 v\!\left(\frac{\zeta(S-P)}{\zeta(H)}\right)
       =v_p\!\left(\frac{a-1}{h}\right).
\]
Here the equality of valuations follows from the formal-logarithm
isometry. The inequality $v_p(h)\ge0$ therefore gives exactly the
source's lower bound for $v_p(a-1)$. In particular the fixed-prime
consequence $|a-1|\ge\ell^{2L+1}$ is retained.

For a comparison extraction, the exact logarithmic ratios are
\[
 \frac{j_i}{h},\qquad\frac{j_i-n_i}{h-1}.
\]
One of $h,h-1$ is a local unit. Therefore one ratio is integral,
and the paired comparison makes the other integral. The fourth-
order logarithmic approximation has the same absolute error as in
the source. With $d_1=1$ and $d_2=d_3=d_4=3$, the paired tests
give, modulo $p^{d_i}$,
\[
 L_i h\equiv j_i,\qquad
 L_i(h-1)\equiv j_i-n_i.
\]
Subtracting yields $L_i\equiv n_i\pmod{p^{d_i}}$.
The four $n_i$ are still integers. Hence the input and Lucas
residuals, their integer divisibilities, and the alternative
$p^3\mid(a-1)$ are precisely those of the original source.
\end{proof}

\begin{theorem}\label{thm:main}
Under the arithmetic premises of the general source, there is a
computable family of integer polynomials in twenty-five variables,
including the input, whose input projection is exactly $\Z$ for
every sufficiently large fixed odd exponent $L$. Their degrees
are at most $4924857429000192<2^{53}$, their support sizes are
strictly below $2^{1193}$, and their integer content is one.
The sufficient exponent remains ineffective.
\end{theorem}
\begin{proof}
Lemma \ref{lem:local} supplies all local divisibilities used in
the source's scalar and aggregate height arguments. Those arguments
involve only the input and the four integer root indices. Their
inequalities and strict exponents are unchanged, so they prove
soundness for every sufficiently large fixed odd $L$.

Every original complete witness remains represented. For each
auxiliary chart $x=g^2$, $y=gk$, substitute the new retained value
$h=g^2$ and new finite root $j=gk$. Since $g$ is a unit, this
identifies the old and new point algebras at all geometric
characters. It preserves the scalar functions, comparisons,
contact functions and pole inputs. The new ordinate and radicand
guards hold on every original prepared character. The retained
pole parameters thus give completeness for every integer input.

The point algebra still has rank $4^9$. Adjoining the two pole
roots gives rank $4^9\cdot4=1048576$. The final pole root may
vanish, while every preceding radicand remains a unit. The same
one-coordinate encoder consequently gives $25$ total variables.
Its degree and support bounds are computed below. Setting the
encoder coordinate to zero gives one, so the integer content is one.
\end{proof}

The proof of Lemma \ref{lem:local} also applies to the corresponding
two actual-trace tests of \cite{Trace}: their local implications supply the same
comparison and scalar extractions. Keeping the four integer root
points and allowing the five auxiliary points to range over $E(K)$
therefore preserves soundness of the twenty-one-coordinate source.
All its retained guards, including the nonterminal square-root
guard, remain in force. Completeness still requires simultaneous
rational realization of the two actual square equations; the
twenty-five-variable pole witnesses do not imply that realization.

\section{Finite specification and verification}
Use the point basis weights and coefficient scales proved in
\cite{FullAux}. The four root blocks have weights $(0,2,2,4)$;
the five auxiliary blocks have weights $(0,2,2,3)$. The diagonal
coefficient scales are $(1,16,16,80)$ in every block. The point
quadratic stages have degrees $(8,12)$ for each root block and
$(6,9)$ for each auxiliary block.

The subalgebra norm and terminal-root construction gives
\[
 \deg G_*\le4696539189,\qquad \deg\Omega\le170852.
\]
Therefore
\[
 \deg\det(I-E_0G_*M_\Omega)
 \le1048576(4696539189+170852+1)
 =4924857429000192.
\]
The native integer circuit has $6169$ arithmetic nodes, $101$
recorded guards and $128$ retained norm factors. The encoder degree
is at most $1048576$. Its support bound is obtained by counting
monomials below the total-degree bound with this encoder-degree
restriction; it is strictly below $2^{1193}$.
For the algebraic fixture $L=3$, the coefficient norm bound is
$2^{34765159307149312}$. No assertion is made that $L=3$ is
a sufficient defining exponent. For each specified odd $L$ the
same finite compiler gives a coefficient bound.

An independent iterative traversal verifies all retained guards,
subalgebra supports, weighted degrees, coefficient bounds and
quadratic-stage envelopes. It rejects deliberate changes to the
degree and guard records. Exact polynomial arithmetic verifies
the point multiplication table, primitive quartic and discriminant.
Independent chord-and-tangent arithmetic verifies the native
detector operations at $72$ finite characters of actual rational
point charts, including half-integral auxiliary indices. These
are finite circuit checks, not a proof of the external arithmetic
inputs or of simultaneous twenty-one-variable realization.

\begin{thebibliography}{9}\raggedright
\bibitem{General} \emph{Twenty-five variables from combined scalar and comparison conditions},
project source, \nolinkurl{papers/few_variables25/few_variables25_merged_20261110_v19.tex}.
\bibitem{Unit} \emph{A single norm determinant for the 25- and 13-variable constructions},
project source, \nolinkurl{papers/finite_algebra/unit_norm_20261114_v28.tex}.
\bibitem{Terminal} \emph{A possibly vanishing terminal quadratic root},
project source, \nolinkurl{papers/finite_algebra/terminal_quadratic_root_20261115_v40.tex}.
\bibitem{FullAux} \emph{Unrestricted auxiliary points in the derivative detector},
project source, \nolinkurl{papers/conditional_few_variables/full_auxiliary11_20261117_v71.tex}.
\bibitem{Trace} \emph{A rational conic criterion using the actual detector trace},
project source, \nolinkurl{papers/pole_parameters/trace_conic_criterion_20261117_v57.tex}.
\end{thebibliography}
\end{document}
